\paragraph{Headline.} Over the full three-year test windows, the tuned
forecast rule is the best policy ($J=0.375$ and 0.365, i.e.\ about 82\%
better than the passive pond). Next come the 2\% fixed opening (0.498,
0.450) and PPO (0.588, 0.483). Year by year, PPO and the rule are
statistically tied ($\Delta J=-0.08\pm0.22$, PPO better in 2 of 6 years).
The three-year metric, the St.~Francis definition, favors the rule because
PPO's overflows cluster in a few years and raise the three-year $Q_p$.
The warm start made a large difference: from scratch, PPO scored
$J\approx1.1$--1.2.

\paragraph{What PPO changed relative to the rule.} PPO started as an exact
copy of the rule and moved toward holding more water. Its outflow is the
oldest of all policies (young fraction 0.10--0.11, versus 0.18--0.20 for
the rule). But it ends both windows at the crest and overflows
6{,}000--12{,}000~m$^3$ per window (the rule: 85--162~m$^3$). Its $Q_p$ is
therefore 2--3 times the rule's. In terms of $J$, the gain in age does not
pay for the loss in peak-flow control.

\paragraph{Why: the training reward is not exactly $J$.} During training
PPO kept improving the reward it was given: the raw reward per hour
improved from $-0.028$ to $-0.0004$. That reward differs from $J$ in two
ways, and both undervalue overflow:
\begin{enumerate}\itemsep1pt
\item the peak term is a smooth surrogate of the exceedance quantile, not
the quantile itself;
\item reward normalization \emph{clips} rewards at $\pm10$ standard
deviations. That stabilized learning, but it also caps the penalty of
exactly the rare, extreme overflow hours that determine $Q_p$.
\end{enumerate}
So PPO optimized the stated reward well, but the stated reward traded
overflow for age at a different rate than $J$ does.

\paragraph{What would likely close the gap.}
\begin{itemize}\itemsep1pt
\item Normalize rewards without clipping (or with a much wider clip), so
the overflow penalty keeps its weight.
\item Select the policy by $J$: evaluate checkpoints during training on
held-out training-period windows and keep the best one, rather than the
last.
\item Keep the policy close to the rule unless it clearly improves on it,
e.g.\ by penalizing deviation from the cloned rule early in training.
\end{itemize}

\paragraph{Lessons.} Two lessons stand out. First, starting from an
engineering rule turned an RL failure into a policy competitive with a
carefully tuned rule. Second, when the true objective (a quantile over
three years) cannot be the reward, every simplification of it, including
the numerical stabilizations, becomes part of the objective the agent
actually optimizes. Whether the result is better must be judged on the
true objective, with tuned baselines, on data the agent has not seen.
